\Needspace{12\baselineskip}
\section{Dated Timeline}

\begin{longtable}{>{\bfseries\RaggedRight\arraybackslash}p{0.18\linewidth}>{\RaggedRight\arraybackslash}p{0.73\linewidth}}
\toprule
\textbf{Date} & \textbf{Event and historical significance}\\
\midrule
\endfirsthead
\toprule
\textbf{Date} & \textbf{Event and historical significance}\\
\midrule
\endhead
\bottomrule
\endfoot
1905--1968 & Marcel Lefebvre's formation, missionary episcopate, African apostolic delegation, conciliar participation, and Spiritan generalate supply the experience and authority behind the later Society.\\
6 June 1969 / 18 August 1970 & Bishop François Charrière permits and then confirms in writing the Fribourg house for seminarians.\\
1 November 1970 & Charrière erects the Society as a diocesan \emph{pia unio}; the same decree approves its statutes for six years \emph{ad experimentum}, with tacit renewal.\\
18 February 1971 & Cardinal John Wright praises the statutes; his letter is favorable testimony, not a pontifical erection decree.\\
11--13 November 1974 & Apostolic visitors inspect Écône. Their complete report is not public in the checked record.\\
21 November 1974 & Lefebvre issues his declaration opposing “eternal Rome” to postconciliar tendencies he rejects.\\
6 May 1975 & Bishop Pierre Mamie withdraws the diocesan approval with Roman support; the Holy See treats suppression as effective, while the Society contests it.\\
29 June / 1, 6, 11, and 22 July 1976 & Lefebvre conducts priestly ordinations despite prohibition; an initial suspension from conferring orders, formal monition, and its receipt precede notification of suspension \emph{a divinis}.\\
11 September / 11 October 1976 & Paul VI meets Lefebvre and then writes a detailed summons to obedience, joining council, liturgy, and authority.\\
1977--1983 & The Society expands; Saint-Nicolas-du-Chardonnet becomes its most visible Paris apostolate; a group of U.S. priests separates in the 1983 internal crisis.\\
1987 & Cardinal Édouard Gagnon conducts an apostolic visitation; the report remains unpublished.\\
5 May 1988 & Lefebvre and Cardinal Ratzinger sign a doctrinal and juridical protocol contemplating a recognized society and an SSPX bishop.\\
17--30 June 1988 & After warnings and failed negotiations over guarantees and timing, Lefebvre and Antônio de Castro Mayer consecrate Fellay, Tissier, Williamson, and de Galarreta without mandate.\\
1--2 July 1988 & The Congregation for Bishops declares penalties; John Paul II issues \emph{Ecclesia Dei adflicta} and creates the \emph{Ecclesia Dei} Commission.\\
25 March 1991 & Lefebvre dies; the Society's government and four bishops enable institutional continuity.\\
24 August 1996 & The Pontifical Council for Legislative Texts explains formal adherence, distinguishing ministerial action and case-specific lay intention from occasional attendance.\\
August 2000 & An SSPX Jubilee pilgrimage to Rome helps reopen sustained contacts; the Society reports more than 5,000 participants.\\
7 July 2007 & Benedict XVI issues \emph{Summorum Pontificum}, widening general access to the 1962 Missal without regularizing the SSPX.\\
21 January / 10 March 2009 & The four surviving bishops' censures are remitted; Benedict distinguishes that personal measure from the Society's lack of status and legitimate ministry.\\
October 2009--April 2011 & Eight CDF--SSPX doctrinal meetings address the Council, tradition, liturgy, ecclesiology, ecumenism, and religious liberty.\\
September 2011--October 2012 & A doctrinal preamble and possible canonical structure are exchanged, but no agreement results; Williamson is expelled from the SSPX in October 2012.\\
2015--2017 & Francis supplies a Jubilee confession faculty, extends it, and establishes a diocesan route for faculties to assist at SSPX marriages.\\
11 July 2018 / 19 January 2019 & Davide Pagliarani becomes superior general; Francis suppresses the \emph{Ecclesia Dei} Commission and transfers its SSPX work to the CDF.\\
2024--2025 & Tissier de Mallerais and the already-expelled Williamson die, leaving two 1988 bishops active in the Society.\\
2 February--29 June 2026 & SSPX announces 1 July consecrations; dialogue, a DDF warning, announcement of four candidates, and Leo XIV's personal appeal do not change the plan.\\
1 July 2026 & De Galarreta, with Fellay as co-consecrator, consecrates Schreiber, Goldade, Poinsinet de Sivry, and Hanappier without mandate.\\
2--3 July 2026 & The DDF issues a penal decree and explanatory note; Pagliarani rejects the penalties and renews the necessity defense.\\
11--16 July 2026 & The Society lodges preliminary recourse against the decree and claims suspensive effect; no checked Holy See response has appeared by the publication cutoff.\\
\end{longtable}
